% =====================================================================
%  Relativistic extension of Chapter IX, Sections 84--86
%  (Non-dissipative Couette flow with axial current, combined fields)
%
%  Based on Chandrasekhar, "Hydrodynamic and Hydromagnetic Stability"
%  Original: §84 -- Stability with current parallel to axis (toroidal B)
%            §85 -- Combined axial + toroidal magnetic fields
%            §86 -- General non-dissipative case
%
%  Relativistic formulation: ideal RMHD in cylindrical geometry;
%  metric signature (-,+,+,+); c kept explicit.
%  See RELATIVISTIC_CONVENTIONS.md for notation and macros.
% =====================================================================
\chapter*{Relativistic Extension of Chapter IX, \S\S\,84--86:
Non-Dissipative Couette Flow with Axial Current and Combined Fields}
\addcontentsline{toc}{chapter}{Relativistic Couette flow with axial current
  and combined fields (\S\S\,84--86)}
\label{ch:rel-ch9-sec84-86}

The classical results of \S\S\,84--86 analyse the stability of
non-dissipative Couette flow in the presence of a toroidal magnetic
field (produced by an axial current) and of combined axial--toroidal
fields.  When the Alfv\'en speed $\vA$ or the rotational velocity
$V = r\Omega$ is no longer negligible compared with $c$, relativistic
corrections to the inertia, the magnetic stresses, and the mode
propagation speeds become essential.  In this chapter we derive the
relativistic generalisations of these stability criteria, always
verifying that perturbation phase speeds remain sub-luminal.

Throughout we work in flat Minkowski spacetime with cylindrical
coordinates $(t,r,\theta,z)$, metric
\begin{equation}
  ds^2 = -c^2\,dt^2 + dr^2 + r^2\,d\theta^2 + dz^2,
  \label{eq:rel9-cyl-metric}
\end{equation}
and adopt the ideal relativistic MHD condition $F^{\mu\nu}u_\nu = 0$
(vanishing electric field in the fluid rest frame).

%======================================================================
\section{Relativistic definitions and equilibrium}
\label{sec:rel9-84-prelim}
%======================================================================

The equilibrium configuration is a differentially rotating fluid column
between coaxial cylinders at $r = R_1$ and $r = R_2$ with angular
velocity $\Omega(r)$, threaded by a toroidal magnetic field
$H_\theta(r)$ (from an axial current) and, in \S\S\,85--86, a uniform
axial field $H_z$.

\subsection{Relativistic enthalpy and magnetic frequencies}

Let $\edensity$ denote the total energy density and $p$ the isotropic
pressure.  The relativistic enthalpy density is
\begin{equation}
  \enthalpy = \edensity + p.
  \label{eq:rel9-w-def}
\end{equation}
The Lorentz factor of the equilibrium flow is
\begin{equation}
  \Lf = \left(1 - \frac{V^2}{c^2}\right)^{-1/2},
  \qquad V = r\Omega.
  \label{eq:rel9-gamma}
\end{equation}

We define the \emph{relativistic Alfv\'en frequency} for the toroidal
field and the \emph{relativistic Alfv\'en frequency} for the axial
field, respectively, as
\begin{equation}
  \Omega_{H,\mathrm{rel}}^{2}
  = \frac{\mu H_\theta^{2}}{4\pi\,r^{2}\bigl(\enthalpy/c^{2}
    + \mu H_\theta^{2}/(4\pi c^{2})\bigr)}
  = \frac{\Omega_H^{2}}{1 + \Omega_H^{2}\,r^{2}/c^{2}
    + \enthalpy/(\rdensity c^{2})_{\mathrm{corr}}},
  \label{eq:rel9-OmH-rel}
\end{equation}
where the classical definition is $\Omega_H^{2}
= \mu H_\theta^{2}/(4\pi\rdensity\,r^{2})$, and
\begin{equation}
  \Omega_{A,\mathrm{rel}}^{2}
  = \frac{\mu H_z^{2}\,k^{2}}{4\pi\bigl(\enthalpy/c^{2}
    + \bsq/(4\pi c^{2})\bigr)}
  = \frac{\Omega_A^{2}}{1 + \vA^{2}/c^{2}
    + \enthalpy/(\rdensity c^{2})_{\mathrm{corr}}},
  \label{eq:rel9-OmA-rel}
\end{equation}
with the classical $\Omega_A^{2} = \mu H_z^{2}k^{2}/(4\pi\rdensity)$.
In the non-relativistic limit $\enthalpy \to \rdensity c^{2}$,
$V/c \to 0$, $\vA/c \to 0$, and both reduce to the classical
frequencies.

More compactly, writing the total relativistic inertial density as
\begin{equation}
  \rho_{\mathrm{eff}}
  = \frac{\Lf^{2}\,\enthalpy}{c^{2}}
    + \frac{\bsq}{4\pi c^{2}},
  \label{eq:rel9-rhoeff}
\end{equation}
where $\bsq = \mu(H_\theta^{2} + H_z^{2})/(4\pi)$ is the
comoving magnetic pressure, the Alfv\'en frequencies become
$\Omega_{H,\mathrm{rel}}^{2} = \mu H_\theta^{2}/(4\pi r^{2}\,
\rho_{\mathrm{eff}})$ and $\Omega_{A,\mathrm{rel}}^{2} =
\mu H_z^{2}k^{2}/(4\pi\,\rho_{\mathrm{eff}})$.

\subsection{Equilibrium force balance}

The relativistic radial equilibrium (from $\nabla_\mu T^{\mu r} = 0$
for the combined fluid+field energy-momentum tensor) reads
\begin{equation}
  \frac{d\Pi_{\mathrm{tot}}}{dr}
  = -\frac{\Lf^{2}\,\enthalpy}{c^{2}}\,\frac{V^{2}}{r}
    + \frac{\mu}{4\pi}\,\frac{H_\theta^{2}}{r},
  \label{eq:rel9-equil}
\end{equation}
where $\Pi_{\mathrm{tot}}$ is the total (gas + magnetic) pressure.
In the non-relativistic limit $\Lf \to 1$,
$\enthalpy/c^{2} \to \rdensity$, and this reduces to the classical
equation~(9-104).

%======================================================================
\section{Relativistic \S\,84: Stability with axial current
  (toroidal $B$-field)}
\label{sec:rel9-84}
%======================================================================

We first treat the case $H_z = 0$, so only the toroidal field
$H_\theta(r)$ is present.  The perturbation ansatz has the same
form as in the classical theory: all perturbed quantities are
proportional to $\exp[i(m\theta + kz + pt)]$.

\subsection{Relativistic perturbation equations}

Linearising the relativistic ideal-MHD equations about the
equilibrium of \S\ref{sec:rel9-84-prelim} and introducing the
Lagrangian displacement $\boldsymbol{\xi}$ via
$\delta u^{i} = i\sigma\,\xi^{i}$ (with $\sigma = p + m\Omega$ the
Doppler-shifted frequency), the relativistic counterparts of the
classical equations~(9-115)--(9-117) are:
\begin{align}
  \biggl[\sigma_{\mathrm{rel}}^{2}
  - m^{2}\Omega_{H,\mathrm{rel}}^{2}
  &- 2r\Omega_{\mathrm{rel}}\frac{d\Omega_{\mathrm{rel}}}{dr}
  + \frac{\mu r}{4\pi\,\rho_{\mathrm{eff}}}
    \frac{d}{dr}\!\left(\frac{H_\theta^{2}}{r^{2}}\right)
  \biggr]\xi_r
  + 2i\bigl(\sigma_{\mathrm{rel}}\,\Omega_{\mathrm{rel}}
    - m\,\Omega_{H,\mathrm{rel}}^{2}\bigr)\xi_\theta
  = D\varpi_{\mathrm{rel}},
  \label{eq:rel9-115}
\end{align}
\begin{equation}
  \bigl(\sigma_{\mathrm{rel}}^{2}
  - m^{2}\Omega_{H,\mathrm{rel}}^{2}\bigr)\xi_\theta
  - 2i\bigl(\sigma_{\mathrm{rel}}\,\Omega_{\mathrm{rel}}
    - m\,\Omega_{H,\mathrm{rel}}^{2}\bigr)\xi_r
  = \frac{im}{r}\,\varpi_{\mathrm{rel}},
  \label{eq:rel9-116}
\end{equation}
\begin{equation}
  \bigl(\sigma_{\mathrm{rel}}^{2}
  - m^{2}\Omega_{H,\mathrm{rel}}^{2}\bigr)\xi_z
  = ik\,\varpi_{\mathrm{rel}},
  \label{eq:rel9-117}
\end{equation}
where $\Omega_{\mathrm{rel}} = \Lf^{2}\Omega$ is the
relativistic effective rotation rate (incorporating the enhanced
inertia), $\sigma_{\mathrm{rel}} = p + m\Omega_{\mathrm{rel}}$,
and $\varpi_{\mathrm{rel}}$ is the relativistic pressure
perturbation normalised by $\rho_{\mathrm{eff}}$.

These equations have \emph{exactly} the same algebraic structure as
the classical set~(9-115)--(9-117), with the replacements
\begin{equation}
  \rdensity \;\longrightarrow\; \rho_{\mathrm{eff}},
  \qquad
  \Omega_H^{2} \;\longrightarrow\; \Omega_{H,\mathrm{rel}}^{2},
  \qquad
  \Omega \;\longrightarrow\; \Omega_{\mathrm{rel}}.
  \label{eq:rel9-dictionary}
\end{equation}

\subsection{The case $m = 0$: Relativistic Michael criterion}
\label{sec:rel9-84-m0}

For $m = 0$ we have $\sigma = p$ and $h_r = h_z = 0$.  The system
reduces to
\begin{equation}
  (DD_{*} - k^{2})\xi_r
  = -\frac{k^{2}}{p^{2}}\,\Psi_{\mathrm{rel}}(r)\,\xi_r,
  \label{eq:rel9-126}
\end{equation}
with the \emph{relativistic Michael discriminant}
\begin{equation}
  \Psi_{\mathrm{rel}}(r)
  = \Phi_{\mathrm{rel}}(r)
    - \frac{\mu}{4\pi\,\rho_{\mathrm{eff}}}\;
      r\frac{d}{dr}\!\left(\frac{H_\theta}{r}\right)^{2},
  \label{eq:rel9-Psi-rel}
\end{equation}
where the relativistic Rayleigh discriminant is
\begin{equation}
  \Phi_{\mathrm{rel}}(r)
  = \frac{1}{r^{3}}\,\frac{d}{dr}\!
    \bigl(\Lf^{2}\,r^{2}\,V\bigr)^{2}
    \;\cdot\frac{1}{\rho_{\mathrm{eff}}}.
  \label{eq:rel9-Phi-rel}
\end{equation}

\begin{relcorrection}
\textbf{Relativistic Michael theorem.}
A necessary and sufficient condition for the stability of the
equilibrium~\eqref{eq:rel9-equil} against axisymmetric perturbations
is that $\Psi_{\mathrm{rel}}(r) > 0$ throughout $(R_1,R_2)$.
The flow is necessarily unstable if $\Psi_{\mathrm{rel}}$
changes sign.

The relativistic correction enters through
(i)~the replacement $\rdensity \to \rho_{\mathrm{eff}}$
(enhanced inertia from enthalpy and magnetic energy),
(ii)~the Lorentz factor $\Lf$ in the centrifugal term, and
(iii)~the modified equilibrium gradient.
In the limit $c \to \infty$,
$\Psi_{\mathrm{rel}} \to \Psi$ as given in the classical
equation~(9-125).
\end{relcorrection}

\subsection{The case $m \ne 0$}
\label{sec:rel9-84-mne0}

For $m \ne 0$, the combined system~\eqref{eq:rel9-115}--\eqref{eq:rel9-117}
with the continuity equation yields, after elimination of
$\xi_\theta$ and $\varpi_{\mathrm{rel}}$,
\begin{align}
  \bigl(\sigma_{\mathrm{rel}}^{2}
  &- m^{2}\Omega_{H,\mathrm{rel}}^{2}\bigr)
  (DD_{*}-k^{2})\xi_r \notag\\
  &= -k^{2}\biggl[\sigma_{\mathrm{rel}}^{2}
    - m^{2}\Omega_{H,\mathrm{rel}}^{2}
    - \mathcal{V}_{\mathrm{rel}}(r)
    + 4\Omega_{\mathrm{rel}}^{2}
    + \frac{4(\sigma_{\mathrm{rel}}\,\Omega_{\mathrm{rel}}
      - m\,\Omega_{H,\mathrm{rel}}^{2})^{2}}
         {\sigma_{\mathrm{rel}}^{2}
           - m^{2}\Omega_{H,\mathrm{rel}}^{2}}
  \biggr]\xi_r,
  \label{eq:rel9-127-rel}
\end{align}
where $\mathcal{V}_{\mathrm{rel}}(r)$ collects the magnetic
shear terms with $\rho_{\mathrm{eff}}$ replacing $\rdensity$.
The integral variational principle~(9-127) carries over verbatim
with all frequencies replaced by their relativistic counterparts
via~\eqref{eq:rel9-dictionary}.

\begin{causalitycheck}
For every mode with Doppler-shifted frequency $\sigma_{\mathrm{rel}}$,
the phase speed must satisfy
\begin{equation}
  v_{\mathrm{ph}}
  = \frac{|\sigma_{\mathrm{rel}}|}{|k|}
  < c.
  \label{eq:rel9-causality-84}
\end{equation}
Since $\sigma_{\mathrm{rel}}^{2}$ is bounded by
$m^{2}\Omega_{H,\mathrm{rel}}^{2} + k^{2}\vA^{2}/(1+\vA^{2}/c^{2})$
and the relativistic Alfv\'en speed
$\vA^{2}/(1+\vA^{2}/c^{2}) < c^{2}$, the causality bound is
automatically satisfied.
\end{causalitycheck}

%======================================================================
\section{Relativistic \S\,85: Combined axial and toroidal fields}
\label{sec:rel9-85}
%======================================================================

We now include a uniform axial magnetic field $H_z$ in addition to
$H_\theta(r)$ and restrict attention to axisymmetric ($m = 0$)
perturbations.

\subsection{Relativistic perturbation equations}

The Lagrangian displacement $\boldsymbol{\xi}$ is defined as in
equations~(9-134) with $p$ replacing $\sigma$.  The induction
relations become
\begin{equation}
  h_r = ikH_z\,\xi_r,
  \qquad
  h_\theta = ikH_z\,\xi_\theta
    - r\xi_r\frac{d}{dr}\!\left(\frac{H_\theta}{r}\right),
  \qquad
  h_z = ikH_z\,\xi_z,
  \label{eq:rel9-135-rel}
\end{equation}
which are identical in form to the classical (9-135) because the
ideal-MHD induction equation is kinematic.

After expressing perturbations in terms of $\boldsymbol{\xi}$
and using the relativistic inertia $\rho_{\mathrm{eff}}$,
the relativistic counterparts of (9-136)--(9-137) read
\begin{equation}
  \bigl(p^{2} - \Omega_{A,\mathrm{rel}}^{2}
    - r\frac{d}{dr}\Omega_{A,\mathrm{rel}}^{2}
    + r\frac{d}{dr}\Omega_{H,\mathrm{rel}}^{2}\bigr)\xi_r
  + 2i\bigl(p\,\Omega_{\mathrm{rel}}
    - \Omega_{A,\mathrm{rel}}\,\Omega_{H,\mathrm{rel}}\bigr)\xi_\theta
  = D\varpi_{\mathrm{rel}},
  \label{eq:rel9-136-rel}
\end{equation}
\begin{equation}
  \bigl(p^{2} - \Omega_{A,\mathrm{rel}}^{2}\bigr)\xi_\theta
  - 2i\bigl(p\,\Omega_{\mathrm{rel}}
    - \Omega_{A,\mathrm{rel}}\,\Omega_{H,\mathrm{rel}}\bigr)\xi_r
  = 0.
  \label{eq:rel9-137-rel}
\end{equation}

Combining with the continuity relation $(p^{2} -
\Omega_{A,\mathrm{rel}}^{2})D_{*}\xi_r = k^{2}\varpi_{\mathrm{rel}}$,
one arrives at the master equation
\begin{align}
  (p^{2} - \Omega_{A,\mathrm{rel}}^{2})(DD_{*} - k^{2})\xi_r
  = -k^{2}\biggl[
    &\frac{d}{dr}\bigl(\sigma_{\mathrm{rel}}^{2}
      - \Omega_{A,\mathrm{rel}}^{2}\bigr) \notag\\
    &+ \frac{4\bigl(p\,\Omega_{\mathrm{rel}}
        - \Omega_{A,\mathrm{rel}}\,\Omega_{H,\mathrm{rel}}
      \bigr)^{2}}{p^{2} - \Omega_{A,\mathrm{rel}}^{2}}
  \biggr]\xi_r.
  \label{eq:rel9-142-rel}
\end{align}
This is the relativistic generalisation of the classical
equation~(9-142).

\subsection{The case $\Omega = 0$: pure magnetic stability
  (relativistic kink and sausage modes)}
\label{sec:rel9-85-Omega0}

Setting $\Omega = 0$ (and hence $\sigma = p$),
equation~\eqref{eq:rel9-142-rel} simplifies to
\begin{equation}
  \kappa_{\mathrm{rel}}\,(DD_{*} - k^{2})\xi_r
  = -k^{2}\biggl\{
    -\frac{d\Omega_{H,\mathrm{rel}}^{2}}{dr}
    + \frac{4\,\Omega_{A,\mathrm{rel}}^{2}\,
            \Omega_{H,\mathrm{rel}}^{2}}
           {\kappa_{\mathrm{rel}}}
  \biggr\}\xi_r,
  \label{eq:rel9-143-rel}
\end{equation}
where
\begin{equation}
  \kappa_{\mathrm{rel}}
  = p^{2} - \Omega_{A,\mathrm{rel}}^{2}.
  \label{eq:rel9-kappa-rel}
\end{equation}

The structure is identical to the classical equation~(9-143) with
$\Omega_H \to \Omega_{H,\mathrm{rel}}$ and
$\Omega_A \to \Omega_{A,\mathrm{rel}}$.  Hence the theorems of
\S\,81(a) carry over.  In particular:

\begin{relcorrection}
\textbf{Relativistic kink/sausage stability criterion ($\Omega = 0$).}
The characteristic values of $\kappa_{\mathrm{rel}}$ (and therefore
of $p^{2}$) are real.  A necessary and sufficient condition for
stability is (cf.\ classical (9-145)):
\begin{equation}
  I_{1,\mathrm{rel}}\;\Omega_{A,\mathrm{rel}}^{2}
  > k^{2}\!\int_{R_1}^{R_2}
    \!\left(4\,\Omega_{H,\mathrm{rel}}^{2}
    + r\frac{d\Omega_{H,\mathrm{rel}}^{2}}{dr}\right)
    r\,\xi_r^{2}\,dr,
  \label{eq:rel9-145-rel}
\end{equation}
where $I_{1,\mathrm{rel}} = \int_{R_1}^{R_2}\!
r\bigl[(D_{*}\xi_r)^{2} + k^{2}\xi_r^{2}\bigr]\,dr$.

Substituting the definitions~\eqref{eq:rel9-OmH-rel} and
\eqref{eq:rel9-OmA-rel}, this becomes
\begin{equation}
  I_{1,\mathrm{rel}}\;
  \frac{H_z^{2}}{4\pi\,\rho_{\mathrm{eff}}}
  > \int_{R_1}^{R_2}
    \frac{1}{4\pi\,\rho_{\mathrm{eff}}}
    \left(\frac{H_\theta^{2}}{r^{2}}
    + \frac{d}{dr}\!\left(\frac{H_\theta}{r}\right)^{\!2}\right)
    r\,\xi_r^{2}\,dr.
  \label{eq:rel9-146-rel}
\end{equation}

The relativistic correction is the replacement
$\rdensity \to \rho_{\mathrm{eff}}$.  Since
$\rho_{\mathrm{eff}} > \rdensity$, the stabilising effect of
the axial field is \emph{weakened} in the relativistic regime:
a stronger axial field is required to suppress kink/sausage
instabilities.

A sufficient condition for stability remains that $H_\theta(r)$
decreases more rapidly than $1/r$, exactly as in the classical case.
\end{relcorrection}

\begin{causalitycheck}
For $\Omega = 0$, the eigenvalue is $p^{2}
= \kappa_{\mathrm{rel}} + \Omega_{A,\mathrm{rel}}^{2}$.
The phase speed is bounded by
\begin{equation}
  v_{\mathrm{ph}}^{2}
  = \frac{p^{2}}{k^{2}}
  \le \Omega_{A,\mathrm{rel}}^{2}/k^{2}
      + |\kappa_{\mathrm{rel}}|/k^{2}
  < \vA^{2}\!\left(1 + \frac{\vA^{2}}{c^{2}}\right)^{\!-1}
    + c_{\mathrm{fast}}^{2}
  < c^{2},
  \label{eq:rel9-causality-85}
\end{equation}
where $c_{\mathrm{fast}}$ is the fast magnetosonic speed.
Since $c_{\mathrm{fast}}^{2} = \cs^{2} + \vA^{2}
- \cs^{2}\vA^{2}/c^{2} < c^{2}$
(see RELATIVISTIC\_CONVENTIONS.md), all mode speeds lie
strictly below $c$.
\end{causalitycheck}

%======================================================================
\section{Relativistic \S\,86: General non-dissipative case}
\label{sec:rel9-86}
%======================================================================

The most general non-dissipative configuration has $\Omega(r) \ne 0$,
$H_\theta(r) \ne 0$, and $H_z \ne 0$, with perturbations of
arbitrary azimuthal mode number $m$.  The relativistic master equation
is~\eqref{eq:rel9-142-rel} (or its $m \ne 0$ counterpart obtained by
restoring $\sigma_{\mathrm{rel}} = p + m\Omega_{\mathrm{rel}}$).

\subsection{Relativistic eigenvalue problem}

The full eigenvalue problem is: find
$p$ such that~\eqref{eq:rel9-142-rel} (with $p \to
\sigma_{\mathrm{rel}} = p + m\Omega_{\mathrm{rel}}$) admits
a non-trivial solution $\xi_r(r)$ vanishing at $r = R_1$ and
$r = R_2$.

Since the algebraic structure is preserved under the
dictionary~\eqref{eq:rel9-dictionary}, all the qualitative
conclusions of the classical theory carry over:

\begin{enumerate}
\item For $m = 0$, the eigenvalues of $p^{2}$ are real and the
  stability boundary is determined by
  $\Psi_{\mathrm{rel}}(r) > 0$ together with
  the combined-field criterion~\eqref{eq:rel9-145-rel}.

\item For $m \ne 0$, the variational formulation of the classical
  equation~(9-127) generalises directly: in the integral relation
  every $\Omega_H$, $\Omega_A$, $\Omega$ is replaced by
  its relativistic counterpart, and $\rdensity$ by
  $\rho_{\mathrm{eff}}$.

\item The full implications of the general $m \ne 0$ case remain
  unexplored even classically (as noted by Chandrasekhar);
  the same is true relativistically.
\end{enumerate}

\subsection{Summary of relativistic modifications}

\begin{relcorrection}
The transition from classical to relativistic non-dissipative
Couette stability with magnetic fields is effected by three
systematic replacements:
\begin{enumerate}
\item \textbf{Inertia:}
  $\rdensity \;\longrightarrow\; \rho_{\mathrm{eff}}
  = \Lf^{2}\enthalpy/c^{2} + \bsq/(4\pi c^{2})$.
  Both the thermal enthalpy and the magnetic energy contribute
  to the effective inertia.

\item \textbf{Centrifugal term:}
  $\Omega \;\longrightarrow\; \Omega_{\mathrm{rel}} = \Lf^{2}\Omega$.
  The Lorentz factor enhances the effective rotation.

\item \textbf{Alfv\'en frequencies:}
  $\Omega_H^{2},\;\Omega_A^{2} \;\longrightarrow\;
  \Omega_{H,\mathrm{rel}}^{2},\;\Omega_{A,\mathrm{rel}}^{2}$,
  given by~\eqref{eq:rel9-OmH-rel} and~\eqref{eq:rel9-OmA-rel}.
  These are \emph{reduced} relative to the classical values because
  the enhanced inertia $\rho_{\mathrm{eff}} > \rdensity$ appears
  in the denominator.
\end{enumerate}

Physical consequences:
\begin{itemize}
\item The stabilising influence of magnetic fields is
  \emph{weakened} at relativistic Alfv\'en speeds because
  $\Omega_{A,\mathrm{rel}} < \Omega_A$ and
  $\Omega_{H,\mathrm{rel}} < \Omega_H$.
\item The destabilising centrifugal effect is
  \emph{enhanced} by $\Lf^{2}$ factors.
\item Both effects conspire to make relativistic Couette flows
  \emph{less stable} than their Newtonian counterparts for a given
  field strength.
\item All mode phase speeds are bounded by the relativistic
  fast magnetosonic speed $\vf < c$, preserving causality.
\end{itemize}
\end{relcorrection}

\subsection{Astrophysical application: kink modes in neutron star
  magnetospheres}
\label{sec:rel9-86-astro}

The toroidal field instabilities of \S\S\,\ref{sec:rel9-84}--\ref{sec:rel9-86}
are directly relevant to neutron star magnetospheres with
$B_\theta \sim 10^{14}$--$10^{16}$~G \cite{Braithwaite2009,Lasky2015}.
The enhanced inertia $\rho_{\mathrm{eff}} > \rdensity$ weakens the
stabilising magnetic tension, requiring a stronger poloidal field to
suppress kink modes.

\begin{figure}[ht]
\centering
\includegraphics[width=0.95\textwidth]{fig_kink_modes_ns.pdf}
\caption{Kink modes in NS magnetospheres.  Left: growth rate vs.\
  wavenumber for several $\Xi$.  Right: critical axial field for
  stabilisation in the $(\Xi, v_A/c)$ plane.}
\label{fig:kink-modes-ns}
\end{figure}

%--------------------------------------------------------------
\subsection{Kruskal--Shafranov condition for neutron star
  magnetospheres: maximum toroidal-to-poloidal field ratio}
\label{sec:rel9-86-KS-NS}
%--------------------------------------------------------------

We now apply the stability criterion of
\S\ref{sec:rel9-85-Omega0} to derive the maximum
toroidal-to-poloidal magnetic field ratio before the onset of
kink instability ($m = 1$) in a neutron star magnetosphere.
The analysis quantifies the Kruskal--Shafranov (KS) condition
for dipole fields spanning $B_{\mathrm{dip}} = 10^{12}$--$10^{15}$\,G.

\paragraph{Setup.}
We model the NS magnetosphere as a cylindrical flux tube of
radius $R \sim R_{\mathrm{NS}} \approx 10\;\mathrm{km}$, with a
poloidal (axial) field $B_{p}$ and a toroidal field $B_{\varphi}$
generated by magnetospheric currents.  The plasma in the
closed magnetosphere has energy density $\edensity$ dominated
by the magnetic energy: $\edensity \approx B^{2}/(8\pi)$.
For the force-free magnetosphere, the effective inertia is
\begin{equation}
\rho_{\mathrm{eff}}
  = \frac{\enthalpy}{c^{2}} + \frac{B^{2}}{8\pi c^{2}}
  \approx \frac{B^{2}}{4\pi c^{2}},
\label{eq:rel9-rhoeff-NS}
\end{equation}
where we have used $\enthalpy/c^{2} \approx B^{2}/(8\pi c^{2})$
for a magnetically dominated plasma.

\paragraph{The KS condition.}
From the stability criterion~\eqref{eq:rel9-145-rel}, the
$m = 1$ kink mode is stable provided the poloidal Alfv\'en
frequency exceeds the toroidal drive.  For a flux tube of
length $L$ along the poloidal field, the fundamental
axial wavenumber is $k = \pi/L$.  The KS condition for
kink stability is
\begin{equation}
\boxed{\;
\frac{B_{\varphi}}{B_{p}}
  < \frac{\pi R}{L}
    \cdot \frac{1}{\sqrt{1 + v_{\mathrm{A},\varphi}^{2}/c^{2}}},
\;}
\label{eq:rel9-KS-NS}
\end{equation}
where $v_{\mathrm{A},\varphi}^{2} = B_{\varphi}^{2}/
(4\pi\rho_{\mathrm{eff}} + B_{\varphi}^{2}/c^{2})$ is the
relativistic toroidal Alfv\'en speed.  In the force-free limit
$v_{\mathrm{A},\varphi} \to c/\sqrt{2}$, so the relativistic
correction factor is $1/\sqrt{1 + 1/2} = \sqrt{2/3} \approx 0.816$.

For a dipolar field line of footpoint colatitude $\theta$, the
field line length is $L \approx 4R_{\mathrm{NS}}/\sin\theta$
and the flux tube radius at the equator is $R \approx R_{\mathrm{NS}}\sin\theta$.
Thus $\pi R/L \approx \pi\sin^{2}\theta/4$.

\paragraph{Quantitative evaluation.}
The critical ratio $(B_{\varphi}/B_{p})_{\mathrm{crit}}$
depends on the field-line geometry but is independent of the
absolute field strength in the force-free regime.  This is
because both $B_{\varphi}$ and $B_{p}$ scale identically with
$B_{\mathrm{dip}}$, and the stability condition involves only
their ratio.  For a field line at $\theta = 45^{\circ}$:
\begin{equation}
\left(\frac{B_{\varphi}}{B_{p}}\right)_{\mathrm{crit}}
  = \frac{\pi\sin^{2}(45^{\circ})}{4}\cdot\sqrt{\frac{2}{3}}
  = \frac{\pi}{8}\cdot\sqrt{\frac{2}{3}}
  \approx 0.321.
\label{eq:rel9-KS-NS-num}
\end{equation}
For comparison, the classical (non-relativistic) value without
the correction factor would be $\pi/8 \approx 0.393$; the
relativistic reduction is $\sim 18\%$.

Table~\ref{tab:rel-KS-NS} presents the critical ratio for
several footpoint colatitudes and the corresponding maximum
twist angle $\Delta\phi_{\mathrm{max}} = (B_{\varphi}/B_{p})_{\mathrm{crit}}\cdot L/R$.

\begin{table}[htbp]
\centering
\caption{Maximum toroidal-to-poloidal field ratio before kink
instability in a NS magnetosphere, for
$B_{\mathrm{dip}} = 10^{12}$--$10^{15}$\,G (the ratio is
independent of $B_{\mathrm{dip}}$ in the force-free limit).
The relativistic correction factor $\sqrt{2/3}$ reduces all
critical ratios by $\approx 18\%$ compared with the classical
KS condition.}
\label{tab:rel-KS-NS}
\begin{tabular}{ccccc}
\hline
$\theta$ (deg) & $\pi R/L$ &
  $(B_{\varphi}/B_{p})_{\mathrm{crit}}^{\mathrm{cl}}$ &
  $(B_{\varphi}/B_{p})_{\mathrm{crit}}^{\mathrm{rel}}$ &
  $\Delta\phi_{\mathrm{max}}$ (rad)\\
\hline
$30^{\circ}$ & $\pi/16 = 0.196$ & $0.196$ & $0.160$ & $0.503$ \\
$45^{\circ}$ & $\pi/8 = 0.393$ & $0.393$ & $0.321$ & $1.005$ \\
$60^{\circ}$ & $3\pi/16 = 0.589$ & $0.589$ & $0.481$ & $1.005$ \\
$75^{\circ}$ & $0.733$ & $0.733$ & $0.599$ & $0.834$ \\
$90^{\circ}$ & $\pi/4 = 0.785$ & $0.785$ & $0.641$ & $0.641\pi$ \\
\hline
\end{tabular}
\end{table}

\paragraph{Implications for field strengths $B_{\mathrm{dip}} = 10^{12}$--$10^{15}$\,G.}
\begin{enumerate}
\item \textbf{Normal pulsars} ($B_{\mathrm{dip}} \sim 10^{12}$\,G):
  the magnetosphere is nearly force-free and the KS condition
  limits the twist to $\Delta\phi \lesssim 1$\,rad.  Larger twists
  drive kink-unstable $m = 1$ modes that dissipate magnetic
  energy---a possible mechanism for magnetospheric activity in
  intermittent pulsars~\citep{Parfrey2013}.

\item \textbf{High-$B$ pulsars} ($B_{\mathrm{dip}} \sim 10^{13}$\,G):
  the same critical ratio applies; the absolute toroidal field at
  kink onset is $B_{\varphi,\mathrm{crit}} \sim 3\times10^{12}$\,G
  (at $\theta = 45^{\circ}$), comparable to the spindown field.

\item \textbf{Magnetars} ($B_{\mathrm{dip}} \sim 10^{14}$--$10^{15}$\,G):
  the kink instability threshold at $B_{\varphi}/B_{p} \approx 0.32$
  implies $B_{\varphi,\mathrm{crit}} \sim 3\times10^{13}$--$3\times10^{14}$\,G.
  Observational evidence from magnetar outbursts suggests that the
  magnetosphere is driven to twists $\Delta\phi \sim 1$--$2$\,rad
  before violent energy release~\citep{Thompson2002,Beloborodov2009},
  consistent with this KS threshold.

\item \textbf{Relativistic vs.\ classical.}
  The $18\%$ reduction in the critical ratio from the relativistic
  correction means that the magnetosphere becomes kink-unstable at
  a lower twist than classically predicted.  For magnetars, this
  translates to a $\sim 18\%$ lower energy threshold for the onset
  of giant flares, which may help reconcile theoretical predictions
  with the observed flare energetics.
\end{enumerate}

\begin{relcorrection}
The Kruskal--Shafranov condition for a NS magnetosphere is
$(B_{\varphi}/B_{p})_{\mathrm{crit}} = (\pi R/L)\sqrt{2/3}$,
where the factor $\sqrt{2/3}$ is the relativistic correction
arising from the enhanced inertia $\rho_{\mathrm{eff}} \approx
B^{2}/(4\pi c^{2})$ in the force-free magnetosphere.  This
reduces the classical critical ratio by $\approx 18\%$ and is
independent of the absolute dipole field strength.
\end{relcorrection}

\medskip
\noindent\textbf{References for \S\S\,84--86.}
Classical: Chandrasekhar~\cite{Chandrasekhar1961}.
NS magnetospheres: Braithwaite~\cite{Braithwaite2009},
Lasky~\cite{Lasky2015}.
Relativistic MHD: Rezzolla \& Zanotti~\cite{RezzollaZanotti2013}.

\subsection{Non-relativistic limit}

Setting $V/c \to 0$ and $\vA/c \to 0$ gives $\Lf \to 1$,
$\rho_{\mathrm{eff}} \to \rdensity$,
$\Omega_{\mathrm{rel}} \to \Omega$,
$\Omega_{H,\mathrm{rel}} \to \Omega_H$, and
$\Omega_{A,\mathrm{rel}} \to \Omega_A$.  Every equation
in this chapter then reduces identically to the corresponding
equation in the classical \S\S\,84--86.

\begin{causalitycheck}
\textbf{Global causality summary.}
In the general non-dissipative case, the perturbation eigenfrequencies
satisfy $|\sigma_{\mathrm{rel}}|/k \le \vf < c$ for all modes.
This follows from the hyperbolicity of the relativistic ideal-MHD
system, whose characteristic speeds are the slow magnetosonic,
Alfv\'en, and fast magnetosonic speeds---all of which lie below $c$
by construction (see the characteristic speed formulae in
RELATIVISTIC\_CONVENTIONS.md):
\begin{equation}
  \vs^{2}
  = \frac{\cs^{2}\,\vA^{2}}{\vf^{2}},
  \qquad
  \vf^{2}
  = \cs^{2} + \vA^{2}
    - \frac{\cs^{2}\,\vA^{2}}{c^{2}}
  < c^{2}.
  \label{eq:rel9-char-speeds}
\end{equation}
No mode in the spectrum can propagate super-luminally.
\end{causalitycheck}
